\chapter{The homotopy category of a cofibration category}
\label{chap:homotcat}

Our goal in this chapter is to describe the homotopy category of a cofibration category. All the definitions and the results of this chapter actually only require the smaller set of precofibration category axioms CF1-CF4. 

For a precofibration category $({\cat M}$, ${\cat W}$, ${\cat Cof})$, recall that we have denoted ${\cat M}_{cof}$ the full subcategory of cofibrant objects of ${\cat M}$. We will show that the functor ${\bf ho}{\cat M}_{cof} \ra {\bf ho} {\cat M}$ is an equivalence of categories (Anderson, \cite{Anderson1}). In fact, we will develop an axiomatic description of {\it cofibrant approximation} functors $t \colon {\cat M}^{'} \ra {\cat M}$, which are modelled on the properties of the inclusion ${\cat M}_{cof} \ra {\cat M}$. Cofibrant approximation functors are left approximations in the sense of \definitionref{defn:wleftapproximation}, and by the Approximation \theoremref{thm:approxthm} they induce an isomorphism at the level of the homotopy category. 

Cofibrant aproximation functors will resurface later in \sectionref{sec:colimarbitrarycat}, when we will reduce the construction of homotopy colimits indexed by arbitrary small diagrams to the construction of homotopy colimits indexed by small direct categories.

We then turn to the study of homotopic maps in a precofibration category. In ${\cat M}_{cof}$ we will define the {\it left} homotopy relation $\simeq_l$ on maps, and show that $f \simeq_l g$ iff $f \simeq g$. The localization of ${\cat M}_{cof}$ modulo homotopy is denoted $\pi {\cat M}_{cof}$. 

We show that the class of weak equivalences between cofibrant objects admits a calculus of fractions in $\pi {\cat M}_{cof}$. As a consequence we obtain a description of ${\bf ho} {\cat M}_{cof}$ whereby any map in ${\bf ho} {\cat M}_{cof}$ can be written (up to homotopy!) as a 'left fraction' $f t^{-1}$, with $t$ a weak equivalence. The theory of homotopic maps and calculus of fractions up to homotopy as described here is due to Brown \cite{Brown}.

\section{Fibrant and cofibrant approximations}
\label{sec:cofapproximations}
We are interested in a class of precofibration category functors which are left approximations (\definitionref{defn:wleftapproximation}), and therefore
\begin{enumerate}
\item
induce isomorphisms when passed to the homotopy category (consequence of the Approximation \theoremref{thm:approxthm}), and 
\item
serve as resolutions for computing total left derived functors (\theoremref{thm:generalexistenceleftderived}). 
\end{enumerate}

These are the {\it cofibrant approximation} functors, defined below. The cofibrant approximation functors should be thought of as an axiomatization of the inclusion ${\cat M}_{cof}$ $\ra$ ${\cat M}$, where ${\cat M}$ is a precofibration category.

\index{approximation!(co)fibrant}
\begin{defn}(Cofibrant approximation functors)
\label{defn::cofapproximation}
\mbox{}
Let ${\cat M}$ be a precofibration category. A functor $t \colon {\cat M}^{'} \ra {\cat M}$ is a {\it cofibrant approximation} if ${\cat M}^{'}$ is a precofibration category with all objects cofibrant and
\begin{description}
\item[CFA1]
$t$ preserves the initial object and cofibrations
\item[CFA2]
A map $f$ of ${\cat M}^{'}$ is a weak equivalence if and only if $tf$ is a weak equivalence
\item[CFA3]
If $A \ra B$, $A \ra C$ are cofibrations of ${\cat M}^{'}$ then the natural map $tB \Sum_{tA} tC \ra t(B \Sum_{A} C)$ is an isomorphism
\item[CFA4]
Any map $f \colon tA \ra B$ factors as $f = r \comp tf^{'}$ with $f^{'}$ a cofibration of ${\cat M}^{'}$ and $r$ a weak equivalence of ${\cat M}$.
\end{description}
\end{defn}

A cofibrant approximation functor in particular sends any object to a cofibrant object, and sends trivial cofibrations to trivial cofibrations. If ${\cat M}$ is a precofibration category, then the inclusion ${\cat M}_{cof} \ra {\cat M}$ is a cofibrant approximation.

The dual definition for prefibration categories is

\begin{defn}(Fibrant approximation functors)
\label{defn::fibapproximation}
\mbox{}
Let ${\cat M}$ be a prefibration category. A functor $t \colon {\cat M}^{'} \ra {\cat M}$ is a {\it fibrant approximation} if ${\cat M}^{'}$ is a prefibration category with all objects fibrant and
\begin{description}
\item[FA1]
$t$ preserves the final object and fibrations
\item[FA2]
A map $f$ of ${\cat M}^{'}$ is a weak equivalence if and only if $tf$ is a weak equivalence
\item[FA3]
If $B \ra A$, $C \ra A$ are fibrations of ${\cat M}^{'}$ then the natural map $ t(B \times_{A} C) \ra tB \times_{tA} tC$ is an isomorphism
\item[FA4]
Any map $f \colon A \ra tB$ factors as $f = tf^{'} \comp s$ with $f^{'}$ a fibration of ${\cat M}^{'}$ and $s$ a weak equivalence of ${\cat M}$.
\end{description}
\end{defn}

We will need the lemmas below. Recall that two maps $f, g \colon A\ra B$ in a caetgry pair $({\cat M}, {\cat W})$ are homotopic $f \simeq g$ by definition if they have the same image in ${\bf ho}{\cat M}$. The prototypic example of homotopic maps in  a precofibration category are the cylinder structure maps $i_0, i_1$ for any cylinder $IA$ on a cofibrant object $A$
\[
\xymatrix { A \Sum A \ar@{>->}[r]^-{i_0+i_1} & IA \ar[r]^{p}_{\sim} & A }
\]
for $p \eps {\cat W}$ implies that $i_0, i_1$ have the same image in ${\bf ho}{\cat M}$.

\begin{lem}
\label{lem::weakequivcancelation}
\mbox{}
\begin{enumerate}
\item
Let $t \colon {\cat M}^{'} \ra {\cat M}$ be a cofibrant approximation of a precofibration category. For any maps $f, g \colon A \ra B$ of ${\cat M}^{'}$ and weak equivalence $b \colon tB \ra B^{'}$ of ${\cat M}$ with $b \comp tf = b \comp tg$, we have that $f \simeq g$.
\item
Let $t \colon {\cat M}^{'} \ra {\cat M}$ be a fibrant approximation of a prefibration category. For any maps $f, g \colon A \ra B$ of ${\cat M}^{'}$ and weak equivalence $a \colon A^{'} \ra tA$ of ${\cat M}$ with $tf \comp a = tg \comp a$, we have that $f \simeq g$.
\end{enumerate}
\end{lem}

\begin{proof}
We only prove (1). We may assume that $f + g \colon A \Sum A \ra B$ is a cofibration. Indeed, for general $f, g$ we factor $f + g$ as a cofibration $f^{'} + g^{'}$ followed by a weak equivalence $r$. The map $tr$ is a weak equivalence, and so is $b \comp tr$. If we proved that $f^{'} \simeq g^{'}$ then it would follow that $f \simeq g$.

So assume that $f + g \colon A \Sum A \ra B$ is a cofibration. Pick a cylinder $IA$, and construct step by step the commutative diagram below
\begin{center}
$\xymatrix { 
    tA \Sum tA \ar@{>->}[r]^-{ti_0 + ti_1} \ar@{>->}[d]_{tf + tg} & tIA \ar[rr]^{tp}_\sim \ar@{>->}[d]^{th} & & A \ar[d]^{b \comp tf = b \comp tg} \\
    tB \ar@{>->}[r]^{tb_1} & tB_1 \ar@{>->}[r]^{tb_2} & tB_2 \ar[r]^{b_3}_\sim & B^{'}    
} $
\end{center}
In this diagram, the bottom horizontal composition is $b \colon tB \ra B^{'}$. We construct $B_1 = B \Sum_{A \Sum A} IA$ with component maps $b_1$ and $h$. $tB_1$ is the pushout of the left square of our diagram, and using the pushout property we construct the map $tB_1 \ra B^{'}$. We then construct $b_3 \comp tb_2$ as the CFA4 factorization of $tB_1 \ra B^{'}$. 

The maps $b_1, b_2, h$ are cofibrations. By the 2 out of 3 axiom CF2, the maps $tb_2 \comp tb_1$ and therefore $b_2b_1$ are weak equivalences. Since $i_0 \simeq i_1$, we get $f \simeq g$.
\end{proof}

\begin{lem}
\label{lem:cleavagechoice}
\mbox{}
\begin{enumerate}
\item
Let $t \colon {\cat M}^{'} \ra {\cat M}$ be a cofibrant approximation of a precofibration category. For any commutative diagram
\begin{center}
$\xymatrix { 
    & tA^{'} \ar[dl]_{ti^{'}} \ar[dr]^{ti^{''}} & & \\
    tD^{'} \ar[drrr]^(.3){\sigma^{'}}  & & tD^{''} \ar[dr]^{\sigma^{''}} & \\
    & tB^{'} \ar[ul]^(.3){tj^{'}} \ar[ur]_(.3){tj^{''}} & & B     
} $
\end{center}
with $j^{'}, j^{''}$ weak equivalences of ${\cat M}^{'}$ and $\sigma^{'}, \sigma^{''}$ weak equivalences of ${\cat M}$ we have that $j^{'-1} i^{'} = j^{''-1} i^{''}$ in ${\bf ho}{\cat M}^{'}$
\item
Let $t \colon {\cat M}^{'} \ra {\cat M}$ be a fibrant approximation of a prefibration category. For any commutative diagram
\begin{center}
$\xymatrix { 
    & tA^{'}  & & A \ar[dlll]_(.25){\tau^{'}}  \ar[dl]^(.3){\tau^{''}} \\
    tD^{'} \ar[ur]^(.7){tq^{'}} \ar[dr]_{tp^{'}}  & & tD^{''} \ar[ul]_(.7){tq^{''}} \ar[dl]^{tp^{''}} & \\
    & tB^{'} & &     
} $
\end{center}
with $q^{'}, q^{''}$ weak equivalences of ${\cat M}^{'}$ and $\tau^{'}, \tau^{''}$ weak equivalences of ${\cat M}$ we have that $ p^{'} q^{'-1} = p^{''} q^{''-1} $ in ${\bf ho}{\cat M}^{'}$
\end{enumerate}
\end{lem}

\begin{proof}
We only prove (1). We may assume that $j^{'}, j^{''}$ are trivial cofibrations. Indeed, in the case of $j^{'}$ suppose that $j^{'} = p^{'} j^{'}_1$ is a Brown factorization with $s^{'}$ a right inverse of $p^{'}$. Then we may replace $i^{'}$, $j^{'}$, $\sigma^{'}$ with $s^{'}i^{'}$, $j^{'}_1$ resp. $\sigma^{'} p^{'}$.

Suppose now that $j^{'}, j^{''}$ are trivial cofibrations. Construct the sum $D^{'} \Sum_{B^{'}} D^{''}$ with component maps the trivial cofibrations $h^{'}$ and $h^{''}$. In the diagram below
\begin{center}
$\xymatrix { 
    & tA^{'} \ar@{>->}[dl]_{ti^{'}} \ar@{>->}[dr]^{ti^{''}} & & \\
    tD^{'} \ar@{>->}[r]^(.4){th^{'}}  & t(D^{'} \Sum_{B^{'}} D^{''}) \ar[drr]^(.7){\sigma} & tD^{''} \ar@{>->}[l]_(.4){th^{''}} & \\
    & tB^{'} \ar@{>->}[ul]^(.3){tj^{'}} \ar@{>->}[ur]_(.3){tj^{''}} & & B     
} $
\end{center}
the bottom triangle is a pushout by CFA3. The map $\sigma$ exists by the universal property of the pushout, since $\sigma^{'} \comp tj^{'} = \sigma^{''} \comp tj^{''}$, and is a weak equivalence by the 2 out of 3 axiom. \lemmaref{lem::weakequivcancelation} (1) applied to $h^{'}i^{'}, h^{''}i^{''}, \sigma$ implies that $h^{'}i^{'} \simeq h^{''}i^{''}$ in ${\cat M}^{'}$, and we conclude that $j^{'-1} i^{'} = j^{''-1} i^{''}$ in ${\bf ho}{\cat M}^{'}$.
\end{proof}

\begin{thm}
\label{thm:cofapproxareleftapprox}
\mbox{}
\begin{enumerate}
\item
Cofibrant approximation functors are left approximations.
\item
Fibrant approximation functors are right approximations.
\end{enumerate}
\end{thm}

\begin{proof}
We only prove (1). If $t \colon {\cat M}^{'} \ra {\cat M}$ is a cofibrant approximation of a precofibration category ${\cat M}$, then $t$ sends weak equivalences to weak equivalences by CFA2. Axiom LAP1 folows from CFA4. To prove the axiom LAP2, , use CFA4 to construct a factorization $f + p \colon tA^{'} \Sum tB^{'} \ra B$ as a cofibration $ti^{'} + tj^{'}$ followed by a weak equivalence $\sigma$. Since $\sigma \comp tj^{'} = p$, the maps $tj^{'}$ and therefore $j^{'}$ are weak equivalences. Axiom LAP2 (2) is proved by \lemmaref{lem:cleavagechoice}.
\end{proof}

As corollaries of the Approximation \theoremref{thm:approxthm}, we note:
\begin{thm}[Anderson]
\label{thm::catcofobjequiv}
\mbox{}
\begin{enumerate}
\item
Given a precofibration category ${\cat M}$, the inclusion $i_{\cat M} \colon {\cat M}_{cof} \ra {\cat M}$ induces an equivalence of categories ${\bf ho}i_{\cat M} \colon {\bf ho}{\cat M}_{cof} \ra {\bf ho}{\cat M}$
\item
Given a prefibration category ${\cat M}$, the inclusion $j_{\cat M} \colon {\cat M}_{fib} \ra {\cat M}$ induces an equivalence of categories ${\bf ho}j_{\cat M} \colon {\bf ho}{\cat M}_{fib} \ra {\bf ho}{\cat M}$
\end{enumerate}
\end{thm}

More generally we have:

\begin{thm}
\label{thm::homotopycatequiv}
\mbox{}
\begin{enumerate}
\item
A cofibrant approximation $t \colon {\cat M}^{'} \ra {\cat M}$ of a precofibration category induces an equivalence of categories ${\bf ho}{\cat M}^{'} \ra {\bf ho}{\cat M}$.
\item
A fibrant approximation $t \colon {\cat M}^{'} \ra {\cat M}$ of a prefibration category induces an equivalence of categories ${\bf ho}{\cat M}^{'} \ra {\bf ho}{\cat M}$.
\end{enumerate}
\end{thm}

It should be noted that the last theorem is actually a particular case of an even more general result of Cisinski, for which we refer the reader to \cite{Cisinski2}, 3.12.

\begin{rem}
\label{rem:intermediatecofcat}
If $t \colon {\cat M}^{'} \ra {\cat M}$ is a cofibrant approximation of a precofibration category ${\cat M}$, suppose that $i \colon {\cat M}_1 \hookrightarrow {\cat M}$ is a subcategory that includes the image of $t$, with weak equivalences and cofibrations induced from ${\cat M}$. Then $i$ as well as the corestriction $t_1 \colon {\cat M}^{'} \ra {\cat M}_1$ of $t$ both define cofibrant approximations. By \theoremref{thm::homotopycatequiv} both functors ${\bf ho}i$, ${\bf ho}t_1$ are equivalences of categories.
\end{rem}

\begin{rem}
\label{rem:intermediatecofcat2}
Suppose that $t \colon {\cat M}^{'} \ra {\cat M}$ is a functor between precofibration categories such that $t$ restricted to ${\cat M}^{'}_{cof}$ is a cofibrant approximation. In view of \theoremref{thm::catcofobjequiv} and \theoremref{thm::homotopycatequiv}, it is not hard to see that $t$ induces a composite equivalence of categories ${\bf ho}{\cat M}^{'} \la {\bf ho}{\cat M}^{'}_{cof} \ra {\bf ho}{\cat M}$. The proper way to formulate this is to say that the {\it total left derived} functor of t is an equivalence, which we will prove as \theoremref{thm:existenceleftderivedequiv} in the next section.
\end{rem}

\section{Total derived functors in cofibration categories}
\label{sec:leftrightderivedcofcat}

The following result describes a sufficient condition for the existence of a total left resp. right derived functor:

\begin{thm}
\label{thm:existenceleftderived}
Let $({\cat M}_2, {\cat W}_2)$ be a category with weak equivalences.
\begin{enumerate}
\item
If ${\cat M}_1$ is a precofibration category and $u \colon {\cat M}_1 \ra {\cat M}_2$ is a functor that sends trivial cofibrations between cofibrant objects to weak equivalences, then $u$ admits a total left derived functor $({\bf L}u, \eps)$. The natural map $\eps \colon ({\bf L}u)(A) \Ra uA$ is an isomorphism for $A$ cofibrant.
\item
If ${\cat M}_1$ is a prefibration category and $u \colon {\cat M}_1 \ra {\cat M}_2$ is a functor that sends trivial fibrations between fibrant objects to weak equivalences, then $u$ admits a total right derived functor $({\bf R}u, \eta)$. The natural map $\eta \colon uA \Ra ({\bf R}u)(A)$ is an isomorphism for $A$ fibrant.
\end{enumerate}
\end{thm}

More generally:
\begin{thm}
\label{thm:existenceleftderivedapprox}
Let $({\cat M}_2, {\cat W}_2)$ be a category with weak equivalences.
\begin{enumerate}
\item
If $t \colon {\cat M}^{'}_1 \ra {\cat M}_1$ is a cofibrant approximation of a precofibration category ${\cat M}_1$ and $u \colon {\cat M}_1 \ra {\cat M}_2$ is a functor such that $ut$ sends trivial cofibrations to weak equivalences, then $u$ admits a total left derived functor $({\bf L}u, \eps)$. The natural map $\eps \colon ({\bf L}u)(tA) \Ra utA$ is an isomorphism for objects $A$ of ${\cat M}^{'}_1$.
\item
If $t \colon {\cat M}^{'}_1 \ra {\cat M}_1$ is a fibrant approximation of a prefibration category ${\cat M}_1$ and $u \colon {\cat M}_1 \ra {\cat M}_2$ is a functor that sends trivial fibrations to weak equivalences, then $u$ admits a total right derived functor $({\bf R}u, \eta)$. The natural map $\eta \colon utA \Ra ({\bf R}u)(tA)$ is an isomorphism for objects $A$ of ${\cat M}^{'}_1$.
\end{enumerate}
\end{thm}

\begin{proof}
We only prove (1). We may assume that ${\cat W}_2$ is saturated. The composition $ut$ sends trivial cofibrations to weak equivalences, therefore using the Brown Factorization \lemmaref{lem:brownfact} $ut$ sends weak equivalences to weak equivalences since ${\cat W}_2$ is saturated. The result follows from \theoremref{thm:generalexistenceleftderived} applied to the cofibrant approximation $t$ and the functor $u$.
\end{proof}

And the next result describes a sufficient condition for the total left (resp. right) derived functor to be an equivalence of categories:

\begin{thm}
\label{thm:existenceleftderivedequiv}
\mbox{}
\begin{enumerate}
\item
If $t \colon {\cat M}_1 \ra {\cat M}_2$ is a functor between precofibration categories such that its restriction $({\cat M}_1)_{cof} \ra {\cat M}_2$ is a cofibrant approximation, then $t$ admits a total left derived functor $({\bf L}t, \eps)$ and ${\bf L}t$ is an equivalence of categories. The natural map $\eps \colon ({\bf L}t)(A) \Ra tA$ is an isomorphism for $A$ cofibrant.
\item
If $t \colon {\cat M}_1 \ra {\cat M}_2$ is a functor between prefibration categories such that its restriction $({\cat M}_1)_{fib} \ra {\cat M}_2$ is a fibrant approximation, then $t$ admits a total right derived functor $({\bf R}t, \eps)$ and ${\bf R}t$ is an equivalence of categories. The natural map $\eta \colon tA \Ra ({\bf R}t)(A)$ is an isomorphism $A$ fibrant.
\end{enumerate}
\end{thm}

\begin{proof}
For part (1), denote $i_{{\cat M}_1} \colon ({\cat M}_1)_{cof} \ra {\cat M}_1$ the inclusion. The functors $i_{{\cat M}_1}$, $ti_{{\cat M}_1}$ are cofibrant approximations and induce equivalences of categories ${\bf ho}i_{{\cat M}_1}$, ${\bf ho}(ti_{{\cat M}_1})$ by \theoremref{thm::homotopycatequiv}. 

From \theoremref{thm:existenceleftderivedapprox} applied to the cofibrant approximation $i_{{\cat M}_1}$ followed by $t$ we see that $t$ admits a total left derived functor $({\bf L}t, \eps)$. Furthermore, $\eps \colon ({\bf L}t)(A) \Ra tA$ is an isomorphism for $A$ cofibrant, therefore ${\bf L}t {\bf ho}i_{{\cat M}_1} = {\bf ho}(ti_{{\cat M}_1})$ and ${\bf L}t$ is an equivalence of categories since ${\bf ho}i_{{\cat M}_1}$ and ${\bf ho}(ti_{{\cat M}_1})$ are equivalences.
\end{proof}

We will also introduce in the context of (co)fibration categories the notion of Quillen adjoint functors and of Quillen equivalences.

\index{Quillen!adjunction}
\index{Quillen!equivalence}
\begin{defn}
\label{defn:quillenadjoint}
Consider four functors
\begin{center}
$\xymatrix {{\cat M}^{'}_1 \ar[r]^{t_1} & {\cat M}_1 \ar@<2pt>[r]^{u_1} & {\cat M}_2 \ar@<2pt>[l]^{u_2} & {\cat M}^{'}_2 \ar[l]_{t_2} }$
\end{center}
where 
\begin{enumerate}
\item
$u_1 \dashv u_2$ is an adjoint pair.
\item
$t_1$ is a cofibrant approximation of a precofibration category ${\cat M}_1$ and $t_2$ is a fibrant approximation of a prefibration category ${\cat M}_2$
\item
$u_1t_1$ sends trivial cofibrations to weak equivalences and $u_2t_2$ sends trivial fibrations to weak equivalences
\end{enumerate}

We then say that $u_1, u_2$ is a {\it Quillen adjoint} pair with respect to $t_1$ and $t_2$. If additionally 
\begin{enumerate}
\item[(4)]
For any objects $A^{'} \eps {\cat M}^{'}_1$, $B^{'} \eps {\cat M}^{'}_2$, a map $u_1t_1A^{'} \ra t_2B^{'}$ is a weak equivalence iff its adjoint $t_1A^{'} \ra u_2t_2B^{'}$ is a weak equivalence
\end{enumerate}
we say that $u_1, u_2$ is a {\it Quillen equivalence} pair with respect to $t_1$ and $t_2$.
\end{defn}

If the functors $t_1, t_2$ are implied by the context, we may refer to $u_1, u_2$ as a Quillen pair of adjoint functors (resp. equivalences) without direct reference to $t_1$ and $t_2$.

\begin{thm}[Quillen adjunction]
\label{thm:quillenadjoint}
\mbox{}
\begin{enumerate}
\item
A Quillen adjoint pair $u_1 \dashv u_2$ with respect to $t_1, t_2$
\begin{center}
$\xymatrix {{\cat M}^{'}_1 \ar[r]^{t_1} & {\cat M}_1 \ar@<2pt>[r]^{u_1} & {\cat M}_2 \ar@<2pt>[l]^{u_2} & {\cat M}^{'}_2 \ar[l]_{t_2} }$
\end{center}
induces a pair of adjoint functors ${\bf L}u_1 \dashv {\bf R}u_2$
\begin{center}
${\bf L}u_1 : {\bf ho}{\cat M}_1 \rightleftarrows {\bf ho}{\cat M}_2 : {\bf R}u_2$
\end{center}
\item
If additionally $u_1, u_2$ satisfy (4l) (resp. (4r)) of \theoremref{thm:generalexistencetotalderivedadjoint} with respect to $t_1, t_2$, then ${\bf L}u_1$ (resp. ${\bf R}u_2$) are fully faithful.
\item
If $u_1, u_2$ is a pair of Quillen equivalences with respect to $t_1, t_2$ then
\begin{center}
${\bf L}u_1 : {\bf ho}{\cat M}_1 \rightleftarrows {\bf ho}{\cat M}_2 : {\bf R}u_2$
\end{center}
is a pair of equivalences of categories.
\end{enumerate}
\end{thm}

\begin{proof}
This is a corollary of \theoremref{thm:generalexistencetotalderivedadjoint}.
\end{proof}

We leave it to the reader to formulate the definition of abstract Quillen partially adjoint functors in the context of cofibration and fibration categories, and to state the analogue in this context of \theoremref{thm:generalexistencetotalderivedadjoint2} 

\comment {
The last definition and theorem extend to the case of what we call partially Quillen adjoint functors.

\index{Quillen!partial adjunction}
\index{Quillen!partial equivalence}
\begin{defn}
\label{defn:quillenadjoint2}
Consider four functors $v_1, v_2, t_1, t_2$
\begin{center}
$\xymatrix{
    {\cat M}^{'}_1 \ar[r]^{v_1} \ar[d]_{t_1} & {\cat M}_2 \\
    {\cat M}_1 & {\cat M}^{'}_2 \ar[l]_{v_2} \ar[u]_{t_2} &
    }$
\end{center}
such that:
\begin{enumerate}
\item
$v_1, v_2$ are partially adjoint with respect to $t_1, t_2$
\item
$t_1$ is a cofibrant approximation of a cofibration category ${\cat M}_1$ and $t_2$ is a fibrant approximation of a fibration category ${\cat M}_2$
\item
$v_1$ sends trivial cofibrations to weak equivalences and $v_2$ sends trivial fibrations to weak equivalences
\end{enumerate}

We then say that $v_1, v_2$ is a {\it Quillen partially adjoint} pair with respect to $t_1$ and $t_2$. If additionally 
\begin{enumerate}
\item[(4)]
For any objects $A^{'} \eps {\cat M}^{'}_1$, $B^{'} \eps {\cat M}^{'}_2$, a map $v_1A^{'} \ra t_2B^{'}$ is a weak equivalence iff its adjoint $t_1A^{'} \ra v_2B^{'}$ is a weak equivalence
\end{enumerate}
we say that $v_1, v_2$ is a {\it Quillen partial equivalence} pair with respect to $t_1$ and $t_2$.
\end{defn}

With this definition we have

\begin{thm}[Quillen partial adjunction]
\label{thm:quillenadjoint2}
\mbox{}
\begin{enumerate}
\item
A Quillen partially adjoint pair $v_1, v_2$ with respect to $t_1, t_2$
\begin{center}
$\xymatrix{
    {\cat M}^{'}_1 \ar[r]^{v_1} \ar[d]_{t_1} & {\cat M}_2 \\
    {\cat M}_1 & {\cat M}^{'}_2 \ar[l]_{v_2} \ar[u]_{t_2} &
    }$
\end{center}
induces a partial adjunction of ${\bf ho}v_1$, ${\bf ho}v_2$ with respect to ${\bf ho}t_1$, ${\bf ho}t_2$, or equivalently induces adjoint functors ${\bf V}_1 \dashv {\bf V}_2$
\begin{center}
${\bf V}_1 : {\bf ho}{\cat M}_1 \rightleftarrows {\bf ho}{\cat M}_2 : {\bf V}_2$
\end{center}
where ${\bf s}_1 \colon {\bf ho}{\cat M}_1 \ra {\bf ho}{\cat M}^{'}_1$ denotes any quasi-inverse of ${\bf ho}t_1$, and ${\bf s}_2 \colon {\bf ho}{\cat M}_2 \ra {\bf ho}{\cat M}^{'}_2$ any quasi-inverse of ${\bf ho}t_2$, and ${\bf V}_1 = {\bf ho}(v_1) {\bf s}_1$, resp. ${\bf V}_2 = {\bf ho}(v_2) {\bf s}_2$.
\item
If additionally $v_1, v_2$ satisfy (4l) (resp. (4r)) of \theoremref{thm:generalexistencetotalderivedadjoint2} with respect to $t_1, t_2$, then ${\bf V}_1$ (resp. ${\bf V}_2$) are fully faithful.
\item
If $v_1, v_2$ is a pair of Quillen partial equivalences with respect to $t_1, t_2$ then ${\bf V}_1, {\bf V}_2$ is a pair of equivalences of categories.
\end{enumerate}
\end{thm}

\begin{proof}
This is a corollary of \theoremref{thm:generalexistencetotalderivedadjoint2}.
\end{proof}
} %end comment



\section{Homotopic maps}
\label{chap:homotopicmaps}
\index{homotopic maps}
We now turn to a more detailed study of homotopic maps in precofibration categories.

Start with ${\cat M}$ a precofibration category, and let $f,g \colon A \ra B$ be two maps with $A, B$ cofibrant. A {\it left homotopy} from $f$ to $g$ is a commutative diagram
\begin{equation}
\label{eqn:lefthomotopy}
\xymatrix{
    A \Sum A \ar[r]^-{f + g} \ar@{>->}[d]_{i_0 + i_1} & B \ar@{>->}[d]^b_{\sim} \\
    IA \ar[r]^H & B^{'}
    }
\end{equation}
with $IA$ a cylinder of $A$ and $b$ a trivial cofibration. We thus have that $Hi_0 = bf$ and $Hi_1 = bg$. The map $H$ is called the {\it left homtopy map} between $f$ and $g$, and we say that the left homotopy goes through the cylinder $IA$ and through the trivial cofibration $b$ . We write $f \simeq_l g$ to say that $f, g$ are left homotopic.

A left homotopy $f \simeq_l g$ with $B = B^{'}$ and $b = 1_B$ is called {\it strict}. We should be careful to point out that Brown uses the notation $\simeq$ differently - to denote what we call strict left homotopy.

Clearly if $f \simeq_l g$ then $f \simeq g$. Our goal will be to show that the notions of homotopy and left homotopy coincide (\theoremref{thm:homotmaps}).

Here is the dual setup for a prefibration category ${\cat M}$. Suppose that $f,g \colon B \ra A$ be two maps with $A, B$ fibrant. A {\it right homotopy} $f \simeq_r g$ is a commutative diagram
\begin{equation}
\label{eqn:righthomotopy}
\xymatrix{
    B^{'} \ar[r]^H \ar@{->>}[d]_b^{\sim} & A^I \ar@{->>}[d]^{(p_0,\,p_1)} \\
    B \ar[r]^-{(f,\,g)} & A \times A
    }
\end{equation}
with $A^I$ a path object of $A$ and $b$ a trivial fibration. A {\it strict} right homotopy additionally satisfies $B = B^{'}$ and $b = 1_B$.

\begin{thm}[Brown, \cite{Brown}]
\label{thm:homotmaps}
\mbox{}
\begin{enumerate}
\item
In a precofibration category ${\cat M}$, two maps $f, g \colon A \ra B$ with $A, B$ cofibrant are left homotopic iff they are homotopic.
\item
In a prefibration category ${\cat M}$, two maps $f, g \colon A \ra B$ with $A, B$ fibrant are right homotopic iff they are homotopic. 
\end{enumerate}
\end{thm}

\begin{proof}
We only prove part (1). Given a left homotopy \equationref{eqn:lefthomotopy}, since $i_0 \simeq i_1$ we have $Hi_0 \simeq Hi_1$, therefore $bf \simeq bg$ and so $f \simeq g$ since $b$ is a weak equivalence.

If $f \simeq g$ are homotopic in ${\cat M}$ on the other hand, by \theoremref{thm::catcofobjequiv} we see that they are homotopic inside ${\cat M}_{cof}$. From  \theoremref{thm::descrhomot} (c) below it will follow that $f \simeq_l g$.
\end{proof}

Let us now work our way to complete the proof of \theoremref{thm:homotmaps}.

Suppose that $A$ is a cofibrant object in the precofibration category ${\cat M}$. If $IA$ is a cylinder of $A$, a {\it refinement} of $IA$ consists of a cylinder $I^{'}A$ and a trivial cofibration $j \colon IA \ra I^{'}A$ such that the diagram below commutes
\begin{center}
$\xymatrix{
    & A \Sum A \ar@{>->}[dl]_{i_0 + i_1} \ar@{>->}[dr]^{i^{'}_0 + i^{'}_1} & \\
    IA \ar[dr]_p^\sim \ar@{>->}[rr]^j & & I^{'}A \ar[dl]^{p^{'}}_\sim \\
    & A &
  }$
\end{center}

If $f \simeq_l g \colon A \ra B$ through $I^{'}A$ with homotopy map $H^{'}$, then $H^{'}j$ defines a homotopy through $IA$.

Note that given any two cylinders $IA$, $I^{'}A$, we can construct a common refinement $I^{''}A$ by factoring $IA \Sum_{A \Sum A} I^{'}A \ra A$ as a cofibration $IA \Sum_{A \Sum A} I^{'}A \ra I^{''}A$ followed by a weak equivalence $I^{''}A \ra A$.

This allows us to prove the following lemma:

\begin{lem}
\label{lem::choosecyl}
\mbox{}
\begin{enumerate}
\item
Let ${\cat M}$ be a precofibration category, and $f \simeq_l g \colon A \ra B$ with $A, B$ cofibrant. Let $IA$ be a cylinder of $A$. Then one can construct a homotopy $f \simeq_l g$ through $IA$.
\item
Let ${\cat M}$ be a prefibration category, and $f \simeq_r g \colon B \ra A$ with $A, B$ fibrant. Let $A^I$ be a path object for $A$. Then one can construct a homotopy $f \simeq_r g$ through $A^I$.
\end{enumerate}
\end{lem}

\begin{proof}
We only prove (1). Assume that there exists a homotopy $f \simeq_l g$ through another cylinder $I^{'}A$. Construct a common refinement $I^{''}A$ of $IA$ and $I^{'}A$. To prove (1), it suffices to construct a homotopy through the refinement $I^{''}A$. In the commutative diagram below
\begin{center}
$\xymatrix{
    A \Sum A \ar[r]^{f + g} \ar@{>->}[d]_{i^{'}_0 + i^{'}_1} & B \ar@{>->}[d]^{b}_{\sim} \\
    I^{'}A \ar[r]^{H^{'}} \ar@{>->}[d]_j^\sim & B^{'} \ar@{>->}[d]^{b^{'}}_\sim \\
    I^{''}A \ar[r]^{H^{''}} & B^{''}
    }$
\end{center}
$H^{'}, B^{'}$ and $b$ define the homotopy $f \simeq_l g$, and $j$ is the cylinder refinement map from $I^{'}A$ to $I^{''}A$. We construct $B^{''}$ as the pushout of $j$ and $H^{'}$, and we have the desired homotopy $H^{''}, B^{''}$ and $b^{'}b$ through $I^{''}A$.
\end{proof}

We can now prove
\begin{thm}
\label{thm::homotopyequivrel}
\mbox{}
\begin{enumerate}
\item
If ${\cat M}$ is a precofibration category, then $\simeq_l$ is an equivalence relation in ${\cat M}_{cof}$. Furthermore if $f \simeq_l g \colon A \ra B$ with $A, B$ cofibrant then
\begin{enumerate}
\item
If $h \colon B \ra C$ with $C$ cofibrant then $hf \simeq_l hg$
\item
If $h \colon C \ra A$ with $C$ cofibrant then $fh \simeq_l gh$
\end{enumerate}
\item
If ${\cat M}$ is a prefibration category, then $\simeq_r$ is an equivalence relation in ${\cat M}_{fib}$. Furthermore if $f \simeq_r g \colon A \ra B$ with $A, B$ fibrant then
\begin{enumerate}
\item
If $h \colon B \ra C$ with $C$ fibrant then $hf \simeq_r hg$
\item
If $h \colon C \ra A$ with $C$ fibrant then $fh \simeq_r gh$
\end{enumerate}
\end{enumerate}
\end{thm}

\begin{proof}
We only prove (1). Clearly $\simeq_l$ is symmetric and reflexive. 

To see that $\simeq_l$ is transitive, assume $f \simeq_l g \simeq_l h \colon A \ra B$. From the previous lemma, we may assume that both homotopies go through the same cylinder $IA$. Denote these homotopies $H_1, B_1^{'}, b_1$ and $H_2, B_2^{'}, b_2$. Taking the pushout of $b_1$ and $b_2$, we obtain homotopies $H_1^{'}, B^{'}, b$ and $H_2^{'}, B^{'}, b$. 

Notice that both diagrams below are pushouts
\begin{center}
$\xymatrix{
    A \ar@{>->}[r]_\sim^{i_0} \ar@{>->}[d]^\sim_{i_1} & IA \ar@{>->}[d]^\sim & & A \Sum A \Sum A \,\,\, \ar@{>->}[rr]^-{A \Sum (i_0 + i_1)} \ar@{>->}[d]_-{(i_0 + i_1) \Sum A} & & A \Sum IA \ar@{>->}[d]^-{i_0 \Sum IA} \\
    IA \ar@{>->}[r]^-\sim & \,\,\, IA \Sum_A IA & & IA \Sum A \,\,\, \ar@{>->}[rr]^-{IA \Sum i_1} & & IA \Sum_A IA
    }$
\end{center}
The factorization $\nabla \colon \xymatrix{A \Sum A \ar@{>->}[r]^-{i_0 \Sum i_1} & IA \Sum_A IA \ar[r]^-{p + p}_-\sim & A}$ is a cylinder: the map $p+p$ is a weak equivalence because of the first diagram, and the map $i_0 \Sum i_1$ is a cofibration as seen if we precompose the second diagram with the cofibration $(i_0, i_2) \colon A \Sum A \ra  A \Sum A \Sum A$. The commutative diagram below then defines a homotopy from $f$ to $h$.
\begin{center}
$\xymatrix{
    A \Sum A \ar[r]^{f + h} \ar@{>->}[d]_{i_0 \Sum i_1} & B \ar@{>->}[d]^{b}_{\sim} \\
    IA \Sum_A IA \ar[r]^-{H^{'}_1 + H^{'}_2} & B^{'}
    }$
\end{center}

To prove (a), let $IA, H, B^{'}, b$ define a homotopy $f \simeq_l g$. In the diagram below let $C^{'}$ be the pushout of $b, h$.
\begin{center}
$\xymatrix{
    A \Sum A \ar[r]^{f + g} \ar@{>->}[d]_{i_0 \Sum i_1} & B \ar[r]^h \ar@{>->}[d]^b_\sim & C \ar@{>->}[d]^c_\sim \\
    IA \ar[r]^H & B^{'} \ar[r]^{h^{'}} & C^{'}
    }$
\end{center}
The outer rectangle defines a homotopy $hf \simeq_l hg$.

For (b), use \lemmaref{lem:existencerelcyl} to construct relative cylinders $IC, IA$ along $h$. From \lemmaref{lem::choosecyl}, we can construct a homotopy $f \simeq_l g$ through $IA$. Precomposing this homotopy with $IC \ra IA$ yields a homotopy $fh \simeq_l gh$.
\end{proof}

For a precofibration category ${\cat M}$, if we factor ${\cat M}_{cof}$ modulo $\simeq_l$ we obtain a category $\pi_l {\cat M}_{cof}$, with same objects as ${\cat M}_{cof}$. By \theoremref{thm::homotopyequivrel}, the morphisms of $\pi_l {\cat M}_{cof}$ are given by $Hom_{\pi_l {\cat M}_{cof}}(A,B) = Hom_{{\cat M}_{cof}}(A,B)_{/ \simeq_l}$. We define weak equivalences in $\pi_l {\cat M}_{cof}$ to be homotopy classes of maps that have one (and hence all) representatives weak equivalece maps of ${\cat M}$.

Of course, in view of \theoremref{thm:homotmaps} ultimately $\pi_l {\cat M}_{cof} \cong \pi {\cat M}_{cof}$. 

For a prefibration category, $\pi_r {\cat M}_{fib}$ denotes the factorization of ${\cat M}_{fib}$ modulo $\simeq_r$. Weak equivalences in $\pi_r {\cat M}_{fib}$ are by definition homotopy classes of maps that have one (and hence all) representatives weak equivalence maps of ${\cat M}$.


\section{Homotopy calculus of fractions}
\label{chap:homotopycalcoffractions}

We will show that for a precofibration category ${\cat M}$, the category $\pi_l {\cat M}_{cof}$ admits a calculus of left fractions in the sense of Gabriel-Zisman with respect to weak equivalences. A nice way to phrase this is to say that ${\cat M}_{cof}$ admits a {\it homotopy} calculus of left fractions. Dually, given a prefibration category ${\cat M}$, its category of fibrant objects ${\cat M}_{fib}$ admits a homotopy calculus of right fractions.

We say that a category pair $({\cat M}$, ${\cat W})$ satisfies the 2 out of 3 axiom provided that for any composable morphisms $f, g$ of ${\cat M}$, if two of $f, g, gf$ are in ${\cat W}$ then so is the third. Weak equivalences in a precofibration category ${\cat M}$ satisfy the 2 out of 3 axiom, and so do weak equivalences in $\pi_l {\cat M}_{cof}$. 

For category pairs $({\cat M}$, ${\cat W})$ satisfying the 2 out of 3 axiom, the Gabriel-Zisman calculus of fractions takes a simplified form that we recall below. The general case of calculus of fractions - when weak equivalences do not necessarily satisfy the 2 out of 3 axiom - is described in \cite{Gabriel-Zisman} at pag. 12.

\index{calculus of fractions}
\begin{thm}[Simplified calculus of left fractions]
\label{thm::gabrielzisman}
Suppose that $({\cat M}, {\cat W})$ is a category pair satisfying the following:
\begin{enumerate}
\item[(a)]
The 2 out of 3 axiom
\item[(b)]
Any diagram of solid maps with $a \eps {\cat W}$
\begin{equation}
\label{eqn:gabrielzismanhypothesis1}
\xymatrix{
  A \ar[r] \ar[d]_a^\sim & B \ar@{-->}[d]^b_\sim \\
  A^{'} \ar@{-->}[r] & B^{'}
}
\end{equation}
extends to a commutative diagram with $b \eps {\cat W}$
\item[(c)]
For any maps $A^{'} \overset{s}{\underset{\sim}{\ra}} A \overset{f}{\underset{g}{\dbra}} B$ with $fs = gs$ and $s \eps {\cat W}$, there exists $B \overset{s^{'}}{\underset{\sim}{\ra}} B^{'}$ in ${\cat W}$ with $s^{'}f = s^{'}g$.
\end{enumerate}
Then:
\begin{enumerate}
\item
Each map in $Hom_{{\bf ho}{\cat M}} (A, B)$ can be written as a left fraction $s^{-1}f$
\begin{center}
$\xymatrix{ A \ar[r]^f & B^{'} & B \ar[l]_-s^-\sim
  }$
\end{center}
with $s$ a weak equivalence. 
\item
Two fractions $s^{-1}f$, $t^{-1}g$ are equal in ${\bf ho}{\cat M}$ if and only if there exist weak equivalences $s^{'}, t^{'}$ as in the diagram below
\begin{equation}
\label{eqn:descrhomotopy1}
\xymatrix{ 
    & & B^{'''} & & \\
    & B^{'} \ar[ur]_\sim^{s^{'}} & &  B^{''} \ar[ul]^\sim_{t^{'}} & \\
    A \ar[ur]^f \ar[urrr]^(.8)g & & & & B \ar[ulll]^(.4)s_(.8)\sim \ar[ul]^\sim_t
  }
\end{equation}
so that $s^{'}s = t^{'}t$ and $s^{'}f = t^{'}g$.
\end{enumerate}
If furthermore weak equivalences are left cancellable, in the sense that for any pair of maps $f, g \colon A \ra B$ and weak equivalence $h \colon B \ra B^{'}$ with $hf = hg$ we have $f = g$, then
\begin{enumerate}
\item[(3)]
Two maps $f, g \colon A \ra B$ are equal in ${\bf ho}{\cat M}$ if and only if $f = g$.
\end{enumerate}
\end{thm}

The dual result for right fractions is
\begin{thm}[Simplified calculus of right fractions]
\label{thm::gabrielzisman2}
Suppose that $({\cat M}, {\cat W})$ is a category pair satisfying the folowing:
\begin{enumerate}
\item[(a)]
The 2 out of 3 axiom
\item[(b)]
Any diagram of solid maps with $a \eps {\cat W}$
\begin{equation}
\label{eqn:gabrielzismanhypothesis2}
\xymatrix{
  B^{'} \ar@{-->}[r] \ar@{-->}[d]_b^\sim & A^{'} \ar[d]^a_\sim  \\
  B \ar[r] & A
}
\end{equation}
extends to a commutative diagram with $b \eps {\cat W}$
\item[(c)]
For any maps $A \overset{f}{\underset{g}{\dbra}} B \overset{t}{\underset{\sim}{\ra}} B^{'}$ with $tf = tg$ and $t \eps {\cat W}$, there exists $A^{'} \overset{t^{'}}{\underset{\sim}{\ra}} A$ in ${\cat W}$ with $ft^{'} = gt^{'}$.
\end{enumerate}
Then:
\begin{enumerate}
\item
Each map in $Hom_{{\bf ho}{\cat M}} (A, B)$ can be written as a right fraction $fs^{-1}$
\begin{center}
$\xymatrix{ A & A^{'} \ar[l]_s^\sim \ar[r]^f & B
  }$
\end{center}
with $s$ a weak equivalence. 
\item
Two fractions $fs^{-1}$, $gt^{-1}$ are equal in ${\bf ho}{\cat M}$ if and only if there exist weak equivalences $s^{'}, t^{'}$ as in the diagram below
\begin{equation}
\label{eqn:descrhomotopy2}
\xymatrix{ 
    & & A^{'''} \ar[dl]^\sim_{s^{'}} \ar[dr]_\sim^{t^{'}} & & \\
    & A^{'} \ar[dl]^\sim_s \ar[drrr]_(.6)f & &  A^{''} \ar[dlll]^\sim_(.2)t \ar[dr]^g & \\
    A & & & & B
  }
\end{equation}
so that $ss^{'} = tt^{'}$ and $fs^{'} = gt^{'}$.
\end{enumerate}
If furthermore weak equivalences are right cancellable, in the sense that for any pair of maps $f, g \colon A \ra B$ and weak equivalence $h \colon A^{'} \ra A$ with $fh = gh$ we have $f = g$, then
\begin{enumerate}
\item[(3)]
Two maps $f, g \colon A \ra B$ are equal in ${\bf ho}{\cat M}$ if and only if $f = g$.
\end{enumerate}
\end{thm}

We only need to supply a 

\begin{proof}[Proof of \theoremref{thm::gabrielzisman}]
We construct a category ${\cat C}$ with objects $Ob {\cat M}$, and maps defined in terms of fractions as explained below.

Fix two objects $A$ and $B$. Consider the set of fractions $s^{-1}f$ as in (1). Denote $\sim$ the relation defined by (2) on the set of fractions $s^{-1}f$ from $A$ to $B$. The relation $\sim$ is clearly reflexive and symmetric. 

To see that the relation is transitive, assume ${s_1}^{-1}f_1 \sim {s_2}^{-1}f_2 \sim {s_3}^{-1}f_3$. We get a commutative diagram
\begin{center}
$\xymatrix{ 
    B^{''}_1 \ar[r]^{u_1} & B^{'''} & B^{''}_2 \ar[l]_{u_2} \\
    B^{'}_1 \ar[u]^{t_1} & B^{'}_2 \ar[ul]_{t_2} \ar[ur]^{t_3} & B^{'}_3 \ar[u]^{t_4} \\
    A \ar[u]^{f_1} \ar[ur]^(.4){f_2} \ar[urr]_(.3){f_3} && B \ar[u]_{s_3} \ar[ul]_(.4){s_2} \ar[ull]^(.3){s_1}
  }$
\end{center}
where the weak equivalences $t_1, t_2$ exist since ${s_1}^{-1}f_1 \sim {s_2}^{-1}f_2$, the weak equivalences $t_3, t_4$ exist since ${s_2}^{-1}f_2 \sim {s_3}^{-1}f_3$ and the weak equivalences $u_1, u_2$ exist from the hypothesis (b) applied to $t_2, t_3$. The compositions $u_1t_1, u_2t_4$ satisfy $u_1t_1f_1 = u_2 t_4 f_3$ and $u_1t_1s_1 = u_2 t_4 s_3$ which shows that ${s_1}^{-1}f_1 \sim {s_3}^{-1}f_3$, and we have proved that $\sim$ is transitive.

We let ${\cat C}(A, B)$ denote the set of fractions from $A$ to $B$ modulo the equivalence relation $\sim$. Given three objects $A, B, C$, we define composition ${\cat C}(A, B) \times {\cat C}(B, C) \ra {\cat C}(A, C)$ as follows. Given fractions $s^{-1}f$, $t^{-1}g$
\begin{center}
$\xymatrix{ 
    & & C^{''} & & \\
    & B^{'} \ar[ur]^{g^{'}} & &  C^{'} \ar[ul]^\sim_{t^{'}} & \\
    A \ar[ur]^f & & B \ar[lu]_s^\sim \ar[ru]_g & & C \ar[ul]^\sim_t
  }$
\end{center}
we use hypothesis (\ref{eqn:gabrielzismanhypothesis1}) to construct an object $C^{''}$, a map $g^{'}$ and a weak equivalence $t^{'}$ such that $g^{'}s = t^{'}g$, and then we define $t^{-1}g \comp s^{-1}f$ as $(t^{'}t)^{-1}(g^{'}f)$. The proof that the definition of composition does not depend on the choices involved uses hypotheses (b) and (c) (we leave this verification to the reader). Given an object $A$, the fraction $(1_A)^{-1}1_A$ is an identity element for the composition.

Define a functor $F \colon {\cat M} \ra {\cat C}$, by $F(A) = A$ and by sending $f \colon A \ra B$ to the fraction $F(f) = (1_B)^{-1}f$. It is not hard to see that $F$ is compatible with composition, and that if $s \colon A \ra B$ is a weak equivalence in ${\cat M}$ then $F(s)$ has $s^{-1}1_B$ as an inverse. 

Since $F$ sends weak equivalences to isomorphisms, it descends to a functor $\overline{F} \colon {\bf ho}{\cat M} \ra {\cat C}$, and it is straightforward to check that any other functor ${\bf ho}{\cat M} \ra {\cat C}^{'}$ factors uniquely through $\overline{F}$. It follows that the category ${\cat C}$ we constructed is equivalent to ${\bf ho}{\cat M}$, which ends the proof of (1) and (2).

If $f, g \colon A \ra B$ are equal in ${\bf ho}{\cat M}$, then by (2) there exists a weak equivalence $h \colon B \ra B^{'}$ such that $hf = hg$. If weak equivalences are left cancellable, then $f = g$, and this proves (3).
\end{proof}

Let us show that given a precofibration category ${\cat M}$, the category $\pi_l {\cat M}_{cof}$ satisfies the hypotheses of \theoremref{thm::gabrielzisman} up to homotopy.

\index{calculus of fractions!homotopic calculus of fractions}
\begin{thm}
\label{thm::homotopycalcoffrations}
\mbox{}
\begin{enumerate}
\item
In a precofibration category
\begin{enumerate}
\item
Any full diagram with cofibrant objects and weak equivalence $a$
\begin{center}
$\xymatrix{
  A \ar[r] \ar[d]_a^\sim & B \ar@{-->}[d]^b_\sim \\
  A^{'} \ar@{-->}[r] & B^{'}
}$
\end{center}
extends to a (strictly) homotopy commutative diagram with $b$ a weak equivalence and $B^{'}$ cofibrant
\item
For any $f,g \colon A \ra B$ with $A, B$ cofibrant
\begin{enumerate}
\item
If there is a weak equivalence $a \colon A^{'} \ra A$ with $A^{'}$ cofibrant such that $fa \simeq_l ga$, then $f \simeq_l g$
\item
If there is a weak equivalence $b \colon B \ra B^{'}$ with $B^{'}$ cofibrant such that $bf \simeq_l bg$, then $f \simeq_l g$
\end{enumerate}
\end{enumerate}
\item
In a prefibration category
\begin{enumerate}
\item
Any full diagram with fibrant objects and weak equivalence $a$
\begin{center}
$\xymatrix{
  B^{'} \ar@{-->}[r] \ar@{-->}[d]_b^\sim & A^{'} \ar[d]^a_\sim  \\
  B \ar[r] & A
}$
\end{center}
extends to a (strictly) homotopy commutative diagram with $b$ a weak equivalence and $B^{'}$ fibrant
\item
For any $f,g \colon A \ra B$ with $A, B$ fibrant
\begin{enumerate}
\item
If there is a weak equivalence $a \colon A^{'} \ra A$ with $A^{'}$ fibrant such that $fa \simeq_r ga$, then $f \simeq_r g$
\item
If there is a weak equivalence $b \colon B \ra B^{'}$ with $B^{'}$ fibrant such that $bf \simeq_r bg$, then $f \simeq_r g$
\end{enumerate}
\end{enumerate}
\end{enumerate}
\end{thm}

\begin{proof}
To prove (1) (a), denote $f \colon A \ra B$ and let $IA$ be a cylinder of $A$. The diagram
\begin{center}
$\xymatrix{
  A \ar[r]^f \ar[d]_a^\sim & B \ar[d]^{b} \\
  A^{'} \ar[r] & A^{'} \Sum_A IA \Sum_A B
}$
\end{center}
is strictly homotopy commutative. It remains to show that $b$ is a weak equivalence. 

Denote $F \colon IA \ra IA \Sum_A B$ the map induced by the $1^{st}$ component of the sum. In the pushout diagram
\begin{center}
$ \xymatrix {
        A \Sum A \ar[r]^{1 \Sum f} \ar@{>->}[d]_{i_0 + i_1} &
	A \Sum B \ar@{>->}[d] & \\
	IA \ar[r]^-F &
	IA \Sum_A B
	}$
\end{center}
the right vertical map is a cofibration and therefore $Fi_0 \colon A \ra IA \Sum_A B$ is a cofibration.

The map $b$ factors as $B \ra IA \Sum_A B \ra A^{'} \Sum_A IA \Sum_A B$. The first factor is a pushout of $i_1 \colon A \ra IA$, therefore a trivial cofibration. The second factor is a pushout of the weak equivalence $a$ by the cofibration $Fi_0$, therefore a weak equivalence by excision.

To prove (1) (b) (i), pick (\lemmaref{lem:existencerelcyl}) relative cylinders $IA^{'}, IA$ over $a$. By \lemmaref{lem::choosecyl}, there exists a homotopy $fa \simeq_l ga$ through $IA^{'}$ and through a trivial cofibration $b^{''}$. We get a commutative diagram
\begin{center}
$\xymatrix{
  A^{'} \Sum A^{'} \ar[r]^{a \Sum a}_\sim \ar@{>->}[d]_{i^{'}_0 + i^{'}_1} & A \Sum A \ar@{>->}[d]_{i_0 + i_1} \ar[r]^{f + g} & B \ar@{>->}[d]_\sim^{b^{''}} \\
  IA^{'} \ar[r]^-{h_1}_-\sim & IA^{'} \Sum_{A^{'} \Sum A^{'}} {A \Sum A} \ar[r]^-{h_2} \ar@{>->}[d]^j_\sim & B^{''} \ar@{>->}[d]^{b^{'}}_\sim \\
  & IA \ar[r]^H & B^{'}
}$
\end{center}
where the map $j$ is a cofibration because $IA^{'}, IA$ are relative cylinders. But $j$ is actualy a trivial cofibration. To see that, notice that since $a$ is a weak equivalence and $A, A^{'}$ are cofibrant, $a \Sum a$ is also a weak equivalence and by excision so is $h_1$. The map $IA^{'} \ra IA$ is a weak equivalence since $a$ is, and by the 2 out of 3 Axiom the map $j$ is a weak equivalence.

We define $B^{'}$ as the pushout of $j$ by $h_2$. The map $b^{'}$ is therefore a trivial cofibration. We let $b = b^{'}b^{''}$, and $IA, H, B^{'}, 1_{B^{'}}$ defines a homotopy $bf \simeq_l bg$.

Let us now prove (1) (b) (ii). Pick a homotopy $bf \simeq_l bg$ going through the cylinder $IA$ and through the trivial cofibration $b^{'}$. In the diagram below
\begin{center}
$\xymatrix{
    A \Sum A \ar[r]^-{f+g} \ar@{>->}[d]_{i_0 + i_1} & B \ar[rr]^-{b}_\sim \ar@{>->}[d]^{b_1} & & B^{'} \ar@{>->}[d]_\sim^{b^{'}} \\
    IA \ar[r]^{h_1} & B_1 \ar@{>->}[r]^{h_2} & B_2 \ar[r]^{h_3}_\sim & B^{''}
}$
\end{center}
construct $b_1$ as the pushout of $i_0 + i_1$ and $h_2, h_3$ as the factorization of $B_1 \ra B^{''}$ as a cofibration followed by a weak equivalence. Notice that $h_2b_1$ is a trivial cofibration, and we have constructed a homotopy $f \simeq_l g$ with homotopy map $h_2h_1$.


The proof of (2) is dual and is omitted.

\end{proof}

The Gabriel-Zisman left calculus of fractions applies therefore to the case of $\pi_l {\cat M}_{cof}$, if ${\cat M}$ is a precofibration category.

\begin{thm}[Brown's homotopy calculus of fractions, \cite{Brown}]
\label{thm::descrhomot}
\mbox{}
\begin{enumerate}
\item
Let ${\cat M}$ be a precofibration category, and $A, B$ be two cofibrant objects.
\begin{enumerate}
\item
Each map in $Hom_{{\bf ho}{\cat M}} (A, B)$ can be written as a left fraction $s^{-1}f$
\begin{center}
$\xymatrix{ A \ar[r]^f & B^{'} & B \ar[l]_-s^-\sim
  }$
\end{center}
with $s$ a weak equivalence and $B^{'}$ cofibrant. 
\item
Two such fractions $s^{-1}f$, $t^{-1}g$ are equal in ${\bf ho}{\cat M}$ if and only if there exist weak equivalences $s^{'}, t^{'}$ as in the diagram (\ref{eqn:descrhomotopy1}) with $B^{'''}$ cofibrant so that $s^{'}s \simeq_l t^{'}t$ and $s^{'}f \simeq_l t^{'}g$.
\item
Two maps $f, g \colon A \ra B$ are equal in ${\bf ho}{\cat M}$ if and only if they are homotopic $f \simeq_l g$
\end{enumerate}
\item
Let ${\cat M}$ be a prefibration category, and $A, B$ be two fibrant objects.
\begin{enumerate}
\item
Each map in $Hom_{{\bf ho}{\cat M}} (A, B)$ can be written as a right fraction $fs^{-1}$
\begin{center}
$\xymatrix{ A & A^{'} \ar[l]_s^\sim \ar[r]^f & B
  }$
\end{center}
with $s$ a weak equivalence and $A^{'}$ fibrant. 
\item
Two such fractions $fs^{-1}$, $gt^{-1}$ are equal in ${\bf ho}{\cat M}$ if and only if there exist weak equivalences $s^{'}, t^{'}$ as in the diagram (\ref{eqn:descrhomotopy2}) with $A^{'''}$ fibrant so that $ss^{'} \simeq_r tt^{'}$ and $fs^{'} \simeq_r gt^{'}$.
\item
Two maps $f, g \colon A \ra B$ are equal in ${\bf ho}{\cat M}$ if and only if $f \simeq_r g$.
\end{enumerate}
\end{enumerate}
\end{thm}

\begin{proof}
This is a consequence of \theoremref{thm::catcofobjequiv}, \theoremref{thm::gabrielzisman} and \theoremref{thm::homotopycalcoffrations}. 
\end{proof}

The proof of \theoremref{thm:homotmaps} is at this point complete, since \theoremref{thm::descrhomot} was its last prerequisite. From this point on we can freely write $\simeq$ instead of $\simeq_l$ and $\simeq_r$.

We can also prove a version \theoremref{thm::descrhomot} that describes ${\bf ho}{\cat M}$ in terms of fractions $s^{-1}f$ with $f, f+s$ cofibrations and $s$ a trivial cofibration: 
\begin{thm}
\label{thm::descrhomotcof}
\mbox{}
\begin{enumerate}
\item
Let ${\cat M}$ be a precofibration category, and $A, B$ be two cofibrant objects.
\begin{enumerate}
\item
Each map in $Hom_{{\bf ho}{\cat M}} (A, B)$ can be written as a left fraction $s^{-1}f$
\begin{center}
$\xymatrix{ A \ar@{>->}[r]^f & B^{'} & B \ar@{>->}[l]_-s^-\sim
  }$
\end{center}
with $f$, $f+s$ cofibrations and $s$ a trivial cofibration.
\item
Two fractions as in \theoremref{thm::descrhomot} (1) (a) $s^{-1}f$, $t^{-1}g$ with $s, t$ trivial cofibrations are equal in ${\bf ho}{\cat M}$ if and only if there exist trivial cofibrations $s^{'}, t^{'}$ as in the diagram (\ref{eqn:descrhomotopy1}) such that $s^{'}s = t^{'}t$ and $s^{'}f \simeq t^{'}g$.
\end{enumerate}
\item
Let ${\cat M}$ be a prefibration category, and $A, B$ be two fibrant objects.
\begin{enumerate}
\item
Each map in $Hom_{{\bf ho}{\cat M}} (A, B)$ can be written as a right fraction $fs^{-1}$
\begin{center}
$\xymatrix{ A & A^{'} \ar@{->>}[l]_s^\sim \ar@{->>}[r]^f & B
  }$
\end{center}
with $f$, $(f, s)$ fibrations and $s$ a trivial fibration
\item
Two fractions as in \theoremref{thm::descrhomot} (2) (a) $fs^{-1}$, $gt^{-1}$ with $s, t$ trivial fibrations are equal in ${\bf ho}{\cat M}$ if and only if there exist trivial fibrations $s^{'}, t^{'}$ as in the diagram (\ref{eqn:descrhomotopy2}) such that $ss^{'} = tt^{'}$ and $fs^{'} \simeq gt^{'}$.
\end{enumerate}
\end{enumerate}
\end{thm}

\begin{proof}
We only prove (1). Denote $\sim$ the equivalence relation defined by \theoremref{thm::descrhomot} (1) (b).

To prove (a) it suffices to show that any fraction $s^{-1}f$ with $s$ a weak equivalence is $\sim$ equivalent to a fraction $t^{-1}g$, with $g, g+t$ cofibrations and $t$ a trivial cofibration. Construct the commutative diagram
\begin{center}
$\xymatrix{
    & B^{'} & \\
    A \ar[ur]^f \ar@{>->}[dr]_{i_A} & B^{''} \ar[u]^\sim_{v} & B \ar[ul]_s^\sim \ar@{>->}[dl]^{i_B} \\
    & A \Sum B \ar@{>->}[u]_u &
  }$
\end{center}
where $vu$ is the factorization of $f + s$ as a cofibration followed by a weak equivalence. Define $g = ui_A$ and $t = ui_b$. The maps $g, g+t$ are cofibrations, the map $t$ is a trivial cofibration and the trivial map $v$ yields the desired $\sim$ equivalence between the fractions $s^{-1}f$ and $t^{-1}g$.

To prove (b), in the diagram below
\begin{center}
$\xymatrix{ 
    & & B^{'''} & & \\
    & B^{'} \ar[ur]_\sim^{s^{'}} & &  B^{''} \ar[ul]^\sim_{t^{'}} & \\
    A \ar[ur]^f \ar[urrr]^(.8)g & & & & B \ar[ulll]^(.4)s_(.8)\sim \ar[ul]^\sim_t
  }$
\end{center}
construct $s^{'}, t^{'}$ as the pushouts of $t, s$. We therefore have $s^{'}s = t^{'}t$. Since $s^{'}f$ and $t^{'}g$ are equal in ${\bf ho}{\cat M}$, we also get $s^{'}f \simeq t^{'}g$.
\end{proof}

Going back to the example of topological spaces $Top$ in \sectionref{sec:exampletopspaces}, recall that we have defined weak equivalences in $Top$ to be the homotopy equivalences given by the classic homotopy relation $\simeq$ in $Top$. We can now show that the definition of classic homotopy $\simeq$ in $Top$ is consistent with \definitionref{defn:homotopicmaps}. 

\begin{prop}
\label{prop:homotopyconsistence}
Two maps $f_0, f_1 \colon A \ra B$ in $Top$ have the same image in ${\bf ho}(Top)$ iff there exists a homotopy $h \colon A \times I \ra B$ which equals $f_k$ when restricted to $A \times {k}$, for $k = 0, 1$.
\end{prop}

\begin{proof}
If a homotopy $h$ exists, clearly $f_0, f_1$ have the same image in ${\bf ho}(Top)$. 

To prove the converse, we observe first that $\xymatrix{A \Sum A \ar[r]^{i_0 + i_1} & IA \ar[r]^p_\sim & A}$ is a cylinder with respect to the Hurewicz cofibration structure in $Top$. Using \theoremref{thm:homotmaps} and \lemmaref{lem::choosecyl}, we construct a left homotopy from $f$ to $g$ through the cylinder $IA$, i.e. a commutative diagram
\[
\xymatrix{
    A \Sum A \ar[r]^-{f + g} \ar@{>->}[d]_{i_0 + i_1} & B \ar@{>->}[d]^b_{\sim} \\
    IA \ar[r]^H & B^{'}
    }
\]
where $b$ is a trivial Hurewicz cofibration. By \lemmaref{lem:trivhurcof}, $b$ admits a retract $r$, and $h = rH$ is the desired homotopy.
\end{proof}
